\section*{Цели и задачи работы}
Определение перемещений с помощью теоретических методов и экспериментально; сравнение результатов расчета и опыта. Задачи:
\begin{itemize}
    \item изучить метод начальных параметров и метод Кастилиано для определения перемещений при изгибе;
    \item провести испытания на установках с консольным стержнем и стержнем на двух шарнирных опорах;
    \item экспериментально измерить прогибы стержня при различных нагрузках.
\end{itemize}

\section*{Краткие теоретические сведения}
При плоском поперечном изгибе перемещения точек оси стержня невелики. Допускается считать, что сечения стержня, оставаясь плоскими, поворачиваются и перемещаются перпендикулярно первоначально недеформированной оси.

Определение перемещений аналитически проводится методом начальных параметров или с помощью теоремы Кастилиано.
Уравнение прогибов по методу начальных параметров имеет вид:
\begin{equation}
    w_x = w_0 - \theta_0 x + \sum_k \frac{P_k (x - a_k)^3}{6 E J_y} + \sum_i \frac{M_i (x - b_i)^2}{2 E J_y} + \sum_j \frac{q_j (x - c_j)^4}{24 E J_y}
\end{equation}

Согласно теореме Кастилиано, обобщенное перемещение равно частной производной потенциальной энергии деформации по обобщенной силе:
\begin{equation}
    \delta_{\text{об}} = \frac{\partial U}{\partial P_{\text{об}}}
\end{equation}

\section*{Схемы установок}
Испытания проводятся на двух типах установок в зависимости от заданного варианта. Схемы представлены на \cref{fig:schemes}.

\begin{figure}[hbt]
    \centering
    \includegraphics[width=0.8\textwidth]{figs/schemes.pdf}
    \caption{Схемы нагружения стержней}
    \label{fig:schemes}
\end{figure}

\section*{Инструкция по проведению эксперимента}
\begin{enumerate}
    \item \textit{Подготовка к испытаниям.}
    Измерьте геометрические параметры стержня и установки.
    Для этого с помощью штангенциркуля и линейки определите длину стержня $l$, ширину $b$, высоту $h$, а также расстояния от места приложения нагрузок до опор и до расчетного сечения.
    \textit{Внимание:} Занесите полученные базовые величины в \cref{tab:params}.

    \item \textit{Сборка схемы.}
    Подготовьте установку в соответствии с вашим вариантом (консоль или двухопорный стержень).
    Для этого разместите стержень на опорах и установите индикатор часового типа строго в том сечении, где требуется определить прогиб.
    \textit{Внимание:} Индикатор должен быть закреплен перпендикулярно оси стержня, касание ножки должно быть плотным.

    \item \textit{Предварительное нагружение.}
    Приложите к стержню начальную нагрузку $P_0$.
    Для этого поместите соответствующий груз на подвес.
    \textit{Внимание:} Эта операция строго обязательна для устранения зазоров в кинематических парах и узлах установки. После стабилизации системы обнулите шкалу индикатора.

    \item \textit{Снятие показаний (Опыт 1).}
    Нагружайте стержень заданными силами $P$ через равные интервалы.
    Для этого последовательно добавляйте грузы на подвес, фиксируя показания индикатора на каждом шаге.
    \textit{Внимание:} Снимайте отсчет только после полного затухания колебаний стрелки. Результаты запишите в столбец «Опыт 1» \cref{tab:measurements}.

    \item \textit{Проверка повторяемости (Опыты 2 и 3).}
    Проведите повторные циклы измерений.
    Для этого разгрузите стержень до величины, несколько меньшей максимальной нагрузки, а затем повторите процедуру нагружения еще дважды.
    \textit{Внимание:} Данные заносятся в соответствующие столбцы \cref{tab:measurements}. Разница между отсчетами на одних и тех же ступенях не должна превышать цену деления прибора.
\end{enumerate}

% Соответствие иллюстраций:
% \label{fig:schemes} — Рис. 4.4 из методических материалов (Схемы нагружения стержней: а — консольный, б — на двух опорах).